% Generated draft: requires mathematical and source review.
\documentclass[11pt]{article}
\usepackage[margin=25mm]{geometry}
\usepackage{fontspec}
\setmainfont{Latin Modern Roman}
\usepackage{amsmath,amssymb,mathtools,unicode-math}
\setmathfont{Latin Modern Math}
\usepackage{enumitem,needspace,xurl,hyperref}
\setlength{\parindent}{0pt}
\setlength{\parskip}{5pt}
\setlength{\emergencystretch}{3em}
\newcommand{\N}{\mathbb{N}}
\newcommand{\Z}{\mathbb{Z}}
\newcommand{\Q}{\mathbb{Q}}
\newcommand{\R}{\mathbb{R}}
\newcommand{\C}{\mathbb{C}}
\newcommand{\F}{\mathbb{F}}
\newcommand{\A}{\mathbb{A}}
\newcommand{\PP}{\mathbb{P}}
\newcommand{\E}{\mathbb{E}}
\newcommand{\eps}{\varepsilon}
\newcommand{\dd}{\,\mathrm{d}}
\begin{document}
\section*{UnsolvedMath ID: 30006583}
\textbf{Source problem code:} catalogue entry
\subsection{Definitions and mathematical statement}
A difference field is a field equipped with a distinguished automorphism \(\sigma\). The theory \(\mathrm{ACFA}\) is the model companion of the theory of difference fields. Let \(K\subseteq\mathcal U\models\mathrm{ACFA}\) be a difference subfield, with the automorphism on \(K\) given by restriction. For a variety with a specified field of definition, its \(\sigma\)-twist \(V^\sigma\) is obtained by applying \(\sigma\) to the coefficients of its defining equations. The supplied formulation considers an algebraic variety \(V\) and a rational map \(\phi:V\dashrightarrow V^\sigma\) over \(K\). A solution is a point \(a\in V(\mathcal U)\) in the domain of definition of \(\phi\) satisfying
\[
\sigma(a)=\phi(a),
\]
where \(\sigma\) acts coordinatewise. Write \(p=\operatorname{tp}(a/K)\) for the complete type of \(a\) over \(K\).

\textbf{Intended conjecture, with unresolved source qualifications.} For every such choice of data and every generic solution \(a\) whose type \(p\) is nonminimal, is
\[
\operatorname{nmdeg}(p)\leq 3
\]
true, with the same constant \(3\) for every dimension \(\dim V\)? Here \(\operatorname{nmdeg}\) denotes the degree of nonminimality named in the supplied record.

The supplied material does not define this invariant or specify the conventions for minimality and genericity in \(\mathrm{ACFA}\). Consequently, the displayed assertion records the intended target but does not supply a complete intrinsic definition of its predicate or invariant. The field of definition of \(V\), any irreducibility or dominance hypotheses, and any additional hypotheses on \(K\), the type, or the characteristic must also be checked against the report. None are silently imposed here. The retained report fragment proposes replacing an earlier expression \(\dim(V)+3\) by \(3\), but omits the preceding assertion and its hypotheses; that earlier bound is therefore not asserted here as a theorem.
\subsection{Short English statement}
The supplied conjecture asks whether every nonminimal generic solution type of an algebraic difference equation has degree of nonminimality at most \(3\), independently of the dimension of the underlying variety. The definitions of the invariant and the precise genericity and hypothesis conventions require source review before this becomes a fully specified mathematical statement.
\subsection{Sources}
\begin{itemize}
\item Supplied local UnsolvedMath registry record, permanent ID \(30006583\), titled \emph{Dimension-Independent Nonminimality of Difference Equations}; exact source problem code \texttt{catalogue entry}. The record belongs to the Oberwolfach Reports Open Problems collection and gives proposed year \(2026\), without identifying a proposer.
\item Supplied report citation: \emph{Model Theory and Arithmetic} (2026), Oberwolfach Reports, DOI 10.4171/owr/2026/3, pp.\ 157--216. Source URL: \url{https://doi.org/10.4171/owr/2026/3}. The supplied \texttt{literature\_sources} entry is labelled ``Oberwolfach Report 3/2026'' and links to \url{https://doi.org/10.4171/owr/2026/3}.
\item Literature assessment dated exactly \texttt{2026-08-22}: the supplied curation labels the problem open and states that the report presents the dimension-independent bound as conjectural and gives no proof. This is supplied dated curation, not an independently verified literature assessment.
\item The registry assigns the statement status \texttt{corrected\_verified} and reports that the official report was checked directly and that OCR, notation, and source-dependent wording were repaired. These are supplied provenance claims; this draft does not independently verify them.
\item The retained original passage mentions a conjecture attributed to Poizat concerning minimality of almost every nonlinear equation, then proposes replacing \(\dim(V)+3\) by \(3\). Its reference ``[2]'' is unresolved in the supplied material. The mention of Poizat does not identify the proposer of the present bound.
\end{itemize}
\end{document}
